%!TEX root = ../main.tex

\section{Discussion}

\subsection{Interpretation of Results}
% PRISMA Item 23a & 23b

\subsubsection{Evolution of the Field, Target Disorders, and Hardware Ecosystems}

The temporal trajectory of publications identified in this review (Figure~\ref{fig:disorders_devices}) demonstrates a marked acceleration following the onset of the COVID-19 pandemic, with prominent publication peaks in 2021 and mid-decade. This pattern mirrors broader macro-trends across digital health and psychiatry, wherein the acute social isolation and healthcare disruptions precipitated by the pandemic triggered both a global surge in psychological distress \cite{santomauro2021covid19depression, chekole2021covidmentalhealth} and an unprecedented demand for scalable, remote, contactless clinical solutions \cite{bufano2023digital}. While early pandemic-era research concentrated heavily on opportunistic sensing and retrospective observation \cite{sanudo2020covidactivity}, the subsequent resurgence reflects a methodological maturation: the field has progressively pivoted from exploratory, descriptive digital phenotyping toward active, context-aware therapeutic deployment.

Regarding diagnostic targets, the literature centers predominantly on internalizing conditions, specifically depressive and anxiety disorders, which constitute the vast majority of included studies. This distribution aligns with contemporary scoping reviews of digital phenotyping \cite{choi2024digitalphenotyping, busshart2026digitalphenotyping, dumas2026digitalphenotyping}. The clinical prioritization of depression and anxiety is grounded in both epidemiological urgency and methodological convenience. Affective and anxious symptom trajectories correspond closely to validated, low-friction digital biomarkers, including reduced spatial exploration through GPS location entropy, altered sleep-wake architecture, diurnal rhythm disruptions, and diminished communication dynamics \cite{saeb2015mobilephone, deangel2022passivemonitoring, shen2025passive}. Furthermore, cognitive behavioral therapy and behavioral activation frameworks for depression and anxiety lend themselves exceptionally well to modular, momentary micro-interventions, such as cognitive reframing prompts or scheduled activity reminders triggered by detected behavioral lapses \cite{balaskas2021ecological, leaning2024smartphone}.

In sharp contrast, severe and episodic psychiatric conditions remain strikingly underrepresented, including neurodevelopmental disorders, bipolar disorder, suicidal ideation, and trauma- and stressor-related disorders. This discrepancy does not reflect a lack of clinical relevance, but rather underscores acute ethical, safety, and algorithmic hurdles \cite{martinezmartin2021ethical}. High-risk conditions such as active suicidal ideation demand fail-safe, real-time crisis escalation protocols; relying on automated sensing triggers introduces substantial clinical liability should false negatives occur \cite{adler2024actionablesensing}. Moreover, modeling complex, multi-phase disorders such as bipolar disorder requires capturing non-linear mood switching and divergent behavioral polarities, such as manic hyper-mobility versus depressive psychomotor retardation, that present formidable predictive challenges for standard mobile algorithms without continuous, high-touch clinical oversight \cite{lanata2015complexity, hassan2025utility}. Meanwhile, emerging applications in sleep-wake, substance use, and somatic symptom disorders reflect targeted domains where behavioral and physiological sensing can be mapped directly to discrete clinical events, such as craving episodes, nocturnal awakenings, or autonomic hyperarousal \cite{hickey2021wearablemental, wang2024contextaware}.

Analysis of hardware platforms in Figure~\ref{fig:disorders_devices} reveals a notable divergence between sensing-only research and intervention delivery. Mobile phones overwhelmingly dominate the research landscape, followed by tablets, whereas wrist-worn wearables and specialized sensors account for only a small fraction. This distribution presents a stark contrast to broader passive-sensing surveys \cite{ometov2021, sheikh2021passivesensing, hickey2021wearablemental}, where smartwatches and fitness bands feature prominently for autonomic and physiological monitoring, including heart rate variability, electrodermal activity, and skin temperature. In context-aware interventions, however, wrist-worn wearables represent only a minor fraction. This discrepancy highlights practical and logistical deployment trade-offs: consumer wearables suffer from rapid battery depletion under continuous high-frequency sensing, proprietary closed operating systems with restrictive background data synchronization rules, fragile Bluetooth tethering pipelines, and high participant non-adherence over multi-week protocols \cite{ometov2021, hassan2025utility, jung2025key}. More critically, the smartphone remains the primary interactive cognitive and therapeutic interface; delivering nuanced, multi-step psychological interventions, such as behavioral activation diaries, guided cognitive reframing, and psychoeducation, directly on miniature smartwatch displays remains ergonomically constrained. Tablets occupy a specialized niche, predominantly deployed in clinician-facilitated sessions, pediatric cohorts, or structured neuropsychological testing tasks requiring expansive visual displays. Finally, specialized sensors, such as standalone Bluetooth beacons, remain restricted to bespoke laboratory prototypes due to procurement costs, maintenance burdens, and lack of consumer scalability.

\subsubsection{Contextual Telemetry: Balancing Granularity, OS Policies, and Privacy}

The Context Mapping matrix (Figure~\ref{fig:bubble_matrix}) articulates how multi-tiered environmental and behavioral indicators are captured and translated into clinical actions. A principal finding of this review is the overwhelming predominance of computational context over physical and social dimensions.

Within computational context, in-app telemetry and micro-temporal patterns represent the core operational signals. This concentration reflects a pragmatic convergence of technical accessibility, operating system constraints, and ecological feasibility. Over recent years, major mobile operating systems (Google Android and Apple iOS) have enacted increasingly stringent background privacy and battery management policies. In particular, aggressive background task termination (such as Android Doze mode and foreground service restrictions in Android 12 to 14) and explicit runtime permissions for continuous background location (such as iOS CoreLocation background tracking prompts and App Tracking Transparency) make sustained, round-the-clock sensor harvesting technically precarious, prone to high data loss, and battery-intensive \cite{dumas2026digitalphenotyping, lee2023digitalphenotyping}. In contrast, in-app telemetry and micro-temporal tracking (session durations, screen-touch latency, module completion cadence, and time-of-day access) execute natively within the application container. They require zero additional sensor permissions, impose negligible battery drain, and run with minimal latency via device-side automation. This enables local, real-time contextual adaptation without requiring continuous, bandwidth-heavy cloud synchronization.

In the physical context domain, mobility and movement indicators dominate, operationalized primarily through GPS geofencing and Bluetooth proximity. This emphasis aligns with foundational digital phenotyping research validating spatial mobility metrics, such as location variance, normalized location entropy, and homestay duration, as robust proxies for depressive psychomotor retardation, social withdrawal, and agoraphobic avoidance \cite{saeb2015mobilephone, onnela2016digitalphenotyping, shen2025passive}. In contrast, ambient environmental metrics (ambient light and acoustic noise) and semantic context are utilized far more sparingly. While continuous acoustic sampling and ambient light monitoring offer valuable insight into domestic routines and circadian exposure \cite{nelson2018smarthome}, they present severe privacy liabilities. Continuous ambient audio recording introduces pervasive eavesdropping concerns and privacy discomfort among clinical participants \cite{martinezmartin2021ethical, sheikh2021passivesensing}, deterring researchers from integrating raw acoustic features into scalable therapeutic pipelines.

In the social context domain, communication metadata and physical proximity sensing provide objective windows into social withdrawal, communication fragmentation, and peer engagement \cite{langener2023capturing}. Intriguingly, clinical and coach support metrics are captured exclusively as contextual background inputs, with zero instances serving as automated tailoring targets. This boundary reflects a critical clinical governance standard: human therapeutic alliance and professional crisis intervention are maintained through structured direct contact or clinical escalation protocols, rather than being delegated to autonomous algorithmic control \cite{rogan2024health, adler2024actionablesensing}.

On the functional target side, adaptive therapy and healthcare manager support emerge as the dominant intervention objectives. This highlights that external context is leveraged predominantly to personalize patient-facing micro-content or to enrich clinician-facing decision-support panels, whereas autonomous learning algorithms remain relatively constrained. This architecture aligns with the paradigm articulated by Adler et al.~\cite{adler2024actionablesensing} and Teles et al.~\cite{teles2017situationawareness}: moving beyond detection in clinical sensing requires closed-loop context models that directly inform both self-management and provider-mediated clinical workflows.

\subsubsection{Intervention Architectures: The Continuum from Deterministic Triggers to Human-in-the-Loop Loops}

Our examination of intervention mechanisms (Figure~\ref{fig:intervention_mechanisms_stacked}) underscores a fundamental paradigm shift away from early, open-loop mobile health designs. Early digital phenotyping was predominantly observational and unidirectional, passively recording telemetry to correlate retrospectively with static clinical rating scales such as the PHQ-9 and GAD-7 \cite{saeb2015mobilephone, onnela2016digitalphenotyping, deangel2022passivemonitoring}. In contrast, the current evidence base demonstrates a pronounced transition toward closed-loop, bidirectional adaptive architectures, operating along a dynamic continuum of active, passive, and hybrid delivery modes.

Triggered delivery represents the primary intervention mechanism across active, passive, and hybrid operationalizations. The prevalence of active triggered interventions reflects a central design principle of Just-In-Time Adaptive Interventions (JITAIs) \cite{nahumshani2018jitai, nahum2025just}: contextual signals are harnessed to identify opportune states of vulnerability and receptivity. Rather than quietly modifying background parameters, systems actively prompt users at contextual decision points to verify internal distress, engage in behavioral activation exercises, or practice cognitive reframing. This empirical distribution offers an important nuance when contrasted with recent systematic reviews \cite{vangenugten2025jitai, van2025beyond, hsu2025personalized}, which noted that mental health JITAIs have historically relied almost exclusively on active, self-reported ecological momentary assessments for intervention tailoring, with passive contextual data rarely operationalized. Our review establishes that when external context is integrated, deterministic rule-based triggers provide an essential, interpretable bridge linking passive context states directly to targeted clinical micro-actions.

Algorithmic mechanisms demonstrate the growing incorporation of machine learning models for risk prediction and dynamic tailoring. However, the limited share of purely passive algorithmic interventions is telling. Clinical researchers remain hesitant to grant machine learning algorithms full autonomy over psychological interventions without user confirmation or clinician validation \cite{adler2024actionablesensing}. Instead, algorithmic outputs are predominantly channeled into hybrid or active configurations that recommend actions while preserving human agency.

This imperative for human oversight is further corroborated by the prominent and balanced distribution of dashboard and manual mechanisms across active, hybrid, and passive configurations. Clinical dashboards bridge automated sensor streams with healthcare workflows, translating raw behavioral metrics into actionable clinical summaries for therapists and care managers \cite{dowding2015dashboards, bosch2025integration}. Similarly, escalation mechanisms embody a crucial safety fail-safe in digital psychiatry: passive continuous monitoring runs in the background until an acute threshold is breached, such as extreme behavioral withdrawal, severe insomnia streaks, or explicit crisis inputs, immediately triggering active client safety plans or alerting clinical supervisors \cite{martinezmartin2021ethical}.

Concrete exemplar architectures identified in our synthesis illustrate these hybrid design paradigms:
\begin{itemize}
    \item \textbf{High-Risk Clinical Cohorts and Hybrid Support:} Fortuna et al.~\cite{fortuna2022peersupportedema} demonstrated the clinical feasibility and acceptability of integrating peer-supported mobile assessment and behavioral tracking for individuals with serious mental illness. By coupling mobile sensing with human peer specialists, the intervention established an essential relational buffer that mitigated digital disengagement and maintained safety in a vulnerable psychiatric population.
    \item \textbf{Actionable Sensing Frameworks:} Adler et al.~\cite{adler2024actionablesensing} demonstrated how multimodal passive sensing must be integrated into clinician-facing dashboards to enable actionable decision support, explicitly addressing the barrier identified by Rogan et al.~\cite{rogan2024health} regarding clinician cognitive overload and algorithmic opacity.
    \item \textbf{Situation-Aware Context Delivery:} Teles et al.~\cite{teles2017situationawareness} operationalized situation-aware architectures that combine spatial trajectories and social metadata to trigger context-tailored cognitive interventions, demonstrating that situating prompts within specific physical and temporal contexts substantially mitigates retrospective recall bias.
    \item \textbf{Contextual Adaptation in Anxiety:} Wang~\cite{wang2024contextaware} formulated personalized, context-aware intervention pipelines targeting social anxiety, coordinating wearable arousal signals with ambient mobile sensing to deliver adaptive in-situ therapeutic prompts during acute social encounters.
    \item \textbf{Passive Behavioral Inference:} Bin Morshed et al.~\cite{binmorshed2021selfesteem} operationalized native Google Activity Recognition API streams to track gross physical movement patterns, demonstrating that background computational abstractions can reliably infer fluctuating psychological states without incurring excessive battery or manual logging costs.
\end{itemize}

Crucially, user-driven interventions were exceedingly rare in our synthesis. This near-total absence demonstrates that contemporary digital mental health systems have largely abandoned static, self-guided repository models in favor of context-aware, proactive triggering.

\subsubsection{Translational Trajectory: The Pilot Treadmill and Methodological Gaps}

The cross-tabulation of study designs and evaluation metrics (Figure~\ref{fig:metrics_designs_heatmap}) provides a clear diagnostic of the field's methodological maturity. Nearly half of all included investigations are pilot and feasibility studies, followed by randomized controlled trials, while formative design studies and observational or quasi-experimental designs represent small minorities.

This heavy concentration within pilot and feasibility frameworks illustrates what implementation scientists term the pilot treadmill or the digital health translation chasm in psychiatry \cite{skivington2021complexinterventions, bosch2025integration}. Deploying context-aware digital interventions requires continuous sensor harvesting, real-time feature engineering, robust mobile operating system compatibility, and complex clinical risk management. Investigators are forced to spend substantial funding cycles establishing basic technical feasibility, battery optimization, and protocol acceptability before advancing to full-scale clinical trials. Furthermore, the rapid obsolescence cycle of mobile operating systems and hardware platforms frequently outpaces the multi-year timeline required to secure funding, obtain institutional ethical approvals, and execute definitive phase-III trials \cite{miralles2020smartphoneapps, leaning2024smartphone}.

Across evaluation dimensions, acceptability and feasibility emerges as a foundational metric across all designs. Similarly, user engagement and adherence is rigorously evaluated in observational trials, randomized controlled trials, and pilot studies. This reflects widespread recognition that digital mental health interventions suffer from steep engagement decay and high participant attrition over time \cite{stone2023evaluation, kraft2024mobile}. Passive sensing directly mitigates the burden of frequent manual logging; however, maintaining sustained participant engagement with triggered therapeutic recommendations remains a critical hurdle.

Importantly, usability and user experience metrics exhibit a precipitous decline as studies advance along the translational continuum: appearing frequently in formative studies and pilot trials, but declining markedly in randomized controlled trials and completely absent from observational designs. This inverse pattern exposes a significant methodological blind spot. As interventions transition to formal clinical efficacy evaluations, researchers frequently treat software interfaces as static artifacts, deprioritizing user-centered design and interaction ergonomics. Yet interaction friction, intrusive notification cadence, and battery drain are primary drivers of trial non-adherence and dropout in digital psychiatry \cite{rogan2024health, stone2023evaluation}. Sustaining clinical efficacy in real-world deployment requires that usability and user experience remain continuous, primary endpoints throughout mature confirmatory trials.

\subsection{Limitations of the Evidence}
% PRISMA Item 23c

The empirical evidence synthesized in this systematic review is characterized by several inherent limitations across clinical, technological, and methodological dimensions:

\begin{enumerate}[label=\textbf{\arabic*.}]
    \item \textbf{Clinical and Intervention Heterogeneity:} The included studies span diverse psychiatric conditions, ranging from mild subclinical depressive symptoms to major depressive disorder, generalized anxiety disorder, and serious mental illness. Concurrently, the deployed therapeutic modalities vary considerably, encompassing cognitive behavioral therapy, behavioral activation, mindfulness, dialectical behavior skills, and crisis stabilization. This substantial clinical and theoretical heterogeneity precludes quantitative meta-analytic pooling and complicates direct comparisons regarding the relative efficacy of specific external context triggers.
    
    \item \textbf{Reliance on Short-Term, Underpowered Pilot Trials:} Nearly half of the evidence base consists of pilot and feasibility trials, which are frequently characterized by small sample sizes and brief intervention windows of only a few weeks. Such short monitoring durations are insufficient to assess long-term clinical efficacy, symptom relapse prevention, longitudinal sensor adherence, or habituation to automated prompts.
    
    \item \textbf{Lack of Standardized Sensing and Feature Protocols:} The field exhibits a pronounced lack of standardized sensing protocols, sampling frequencies, and data preprocessing pipelines. For instance, physical mobility was variously operationalized across studies via raw coordinates, proprietary geofencing boundaries, normalized location entropy, or coarse step counts, with diverse epoch aggregations. This methodological fragmentation hampers cross-study reproducibility, impedes algorithmic benchmark comparisons, and limits the generalizability of predictive models \cite{leaning2024smartphone, deangel2022passivemonitoring}.
    
    \item \textbf{Algorithmic Transparency and Explainability:} Many intervention architectures rely either on rigid, hand-crafted deterministic heuristics (which lack individualized nuance) or proprietary, black-box machine learning models that obscure how specific contextual states translate into clinical recommendations. The lack of algorithmic explainability creates understandable hesitation among clinicians, who remain accountable for patient care yet cannot audit the underlying decision logic \cite{rogan2024health, adler2024actionablesensing}.
    
    \item \textbf{Demographic, Socioeconomic, and Platform Selection Biases:} The majority of included trials recruited predominantly young, tech-literate, urban participants, often drawn from university student pools or online volunteer cohorts. Marginalized demographic groups, socioeconomically disadvantaged individuals, older adults, and populations with severe cognitive or functional impairments remain markedly underrepresented. Furthermore, systematic hardware disparities, such as differential background sensing restrictions across premium devices versus budget smartphones with aggressive vendor-specific battery managers, introduce platform-level biases that distort data completeness and intervention delivery across socioeconomic strata.
\end{enumerate}

\subsection{Limitations of the Review Process}
% PRISMA Item 23d

Several methodological constraints inherent to the systematic review process must also be acknowledged:

\begin{enumerate}[label=\textbf{\arabic*.}]
    \item \textbf{Database Coverage and Language Exclusions:} While our search strategy encompassed four major multidisciplinary and biomedical databases (Scopus, Web of Science, PsycINFO, and PubMed) without publication date limits, the search was explicitly restricted to English-language peer-reviewed articles and conference proceedings. Consequently, relevant empirical investigations published in other languages, regional journals, or technical conference tracks indexed exclusively in specialized engineering repositories may have been omitted.
    
    \item \textbf{Exclusion of Grey Literature and Commercial App Store Ecosystems:} In strict accordance with PRISMA standards for peer-reviewed evidence, our eligibility criteria excluded non-peer-reviewed grey literature, preprints, patents, and commercial mobile application repositories (such as the Apple App Store and Google Play). Numerous commercial mental health applications currently incorporate context-aware push notifications and device-usage tracking; however, because their algorithmic architectures and clinical efficacy outcomes are rarely published in peer-reviewed scientific literature, their technical implementations could not be evaluated.
    
    \item \textbf{Scope Boundary on External versus Internal Context:} To address the distinct "why" gap in ubiquitous computing, our eligibility criteria deliberately restricted inclusion to studies that captured and integrated external contextual information (physical, social, and computing/temporal factors), excluding interventions relying exclusively on internal physiological tracking (such as standalone heart rate variability or electrodermal activity biofeedback) or pure self-report symptom rating scales without external context. While this restriction was necessary to answer our core research questions, it intentionally set aside valuable literature on standalone autonomic biofeedback.
    
    \item \textbf{Publication and Positive Outcome Bias:} In digital health and mobile sensing, there exists a well-documented publication bias favoring positive feasibility, acceptability, and efficacy outcomes. Studies encountering severe technical failures, such as app crashes, massive sensor data loss, unsustainable battery drain, or severe participant dropout, or null clinical outcomes are substantially less likely to be submitted or accepted for publication, potentially skewing the literature toward an overly optimistic appraisal of system reliability.
    
    \item \textbf{Technological Turnover and Publication Latency:} The pace of consumer smart hardware and mobile operating system development outstrips academic publication cycles. The multi-year lag between study design, grant procurement, ethical approval, participant recruitment, trial completion, and peer-reviewed publication means that some reviewed architectures reflect operating system APIs and hardware specifications that have since been modified or deprecated by platform vendors.
\end{enumerate}